\documentclass{amsart}
\usepackage{amsmath,amssymb,amsthm,hyperref}
\newtheorem{theorem}{Theorem}
\newtheorem{lemma}[theorem]{Lemma}
\newtheorem{proposition}[theorem]{Proposition}
\newtheorem{corollary}[theorem]{Corollary}
\theoremstyle{definition}
\newtheorem{remark}[theorem]{Remark}
\DeclareMathOperator{\AI}{AI}
\DeclareMathOperator{\supp}{supp}
\DeclareMathOperator{\RM}{RM}
\DeclareMathOperator{\rank}{rank}
\newcommand{\F}{\mathbb{F}}
\begin{document}

\title{The proportion of balanced Boolean functions with optimal algebraic immunity}
\maketitle

\section{Statement and main results}

Throughout, $n\ge 1$ and $N=2^{n}$. We identify a Boolean function
$f:\F_2^n\to\F_2$ with its truth table $(f(x))_{x\in\F_2^n}\in\F_2^{N}$
and with its support $\supp(f)=\{x: f(x)=1\}$. A function is
\emph{balanced} when $|\supp(f)|=2^{n-1}=N/2$; the set of balanced
functions is $B_n$, so
\begin{equation}\label{eq:Bn}
|B_n|=\binom{N}{N/2}.
\end{equation}
Every $f$ has a unique algebraic normal form
$f(x)=\sum_{u\in\F_2^n}c_u\,x^{u}$ with $x^{u}=\prod_{i:u_i=1}x_i$, and
$\deg f=\max\{\mathrm{wt}(u):c_u\ne0\}$. A nonzero $g$ is an
\emph{annihilator} of $f$ if $fg=0$, equivalently
$\supp(f)\subseteq\{x:g(x)=0\}$. By definition $\AI(f)$ is the least
degree of a nonzero annihilator of $f$ or of $f\oplus 1$, and it is
classical that $\AI(f)\le\lceil n/2\rceil$ for all $f$; we set
\[
m:=\lceil n/2\rceil,\qquad k:=m-1=\lceil n/2\rceil-1 .
\]
$B_n^*$ is the set of balanced $f$ with $\AI(f)=m$, and
$\rho(n)=|B_n^*|/|B_n|$.

\medskip
\noindent\textbf{Summary of what is established.}
We prove unconditionally:
\begin{itemize}
\item an exact matrix reduction characterising $B_n^*$
      (Theorem~\ref{thm:reduction});
\item a parity dichotomy: for even $n$ the relevant matrix $E_S$ is
      strictly tall, for odd $n$ it is square
      (Theorem~\ref{thm:dichotomy});
\item for odd $n$, a complement symmetry reducing $\rho(n)$ to a single
      matrix--invertibility probability (Theorem~\ref{thm:odd-sym});
\item for every odd $n\ge3$, $\,0<\rho(n)<1$, with the lower bound proved
      by an explicit self-contained construction
      (Proposition~\ref{prop:oddbounds});
\item for even $n$, a finite union upper bound for $1-\rho(n)$ whose
      minimum--weight part decays super-exponentially in $n$
      (Proposition~\ref{prop:evenmin}), together with an exact generating
      identity \eqref{eq:genid}.
\end{itemize}
The two limiting statements are genuinely deep and are \emph{not} proved
here in full; we isolate each obstruction precisely. For even $n$ the
limit $\rho(n)\to1$ is reduced to one classical analytic input on the
weight distribution of low-order Reed--Muller codes
(\eqref{eq:cited}); for odd $n$ the value of $\lim\rho(n)$ is left as an
open problem, and we explain why the natural heuristic
$\prod_{i\ge1}(1-2^{-i})$ is unverified.

\begin{theorem}[Matrix reduction]\label{thm:reduction}
Let $E$ be the $N\times M$ matrix over $\F_2$ whose rows are indexed by the
points $x\in\F_2^n$, whose columns are indexed by the monomials $x^{u}$
with $\mathrm{wt}(u)\le k$, and whose $(x,u)$ entry is $x^{u}$, where
\begin{equation}\label{eq:M}
M=\sum_{j=0}^{k}\binom nj .
\end{equation}
For $S\subseteq\F_2^n$ write $E_S$ for the submatrix of rows indexed by
$S$. Then a balanced $f$ lies in $B_n^*$ iff
\[
\rank_{\F_2}E_{\supp(f)}=M\quad\text{and}\quad
\rank_{\F_2}E_{\,\F_2^n\setminus\supp(f)}=M .
\]
\end{theorem}

\begin{theorem}[Parity dichotomy]\label{thm:dichotomy}
With $M$ as in \eqref{eq:M} and $|S|=N/2$:
\begin{enumerate}
\item if $n$ is even then $k=n/2-1$ and $M=N/2-\tfrac12\binom n{n/2}<N/2$
      ($E_S$ is strictly tall);
\item if $n$ is odd then $k=(n-1)/2$ and $M=N/2$ ($E_S$ is square).
\end{enumerate}
\end{theorem}

\begin{theorem}[Odd $n$: complement symmetry]\label{thm:odd-sym}
Let $n$ be odd. For every $S$ with $|S|=N/2$,
\[
\rank_{\F_2}E_S=M\iff \rank_{\F_2}E_{S^c}=M,\qquad S^c=\F_2^n\setminus S.
\]
Hence, for odd $n$, $f\in B_n^*$ iff $E_{\supp(f)}$ is invertible, and
\begin{equation}\label{eq:odd-rho}
\rho(n)=\Pr\bigl[E_{S}\ \text{invertible over }\F_2\bigr],\qquad
S\ \text{uniform of size }N/2 .
\end{equation}
\end{theorem}

\begin{proposition}[Even $n$: the union bound and its dominant term]
\label{prop:evenmin}
For even $n$, with $w_0:=2^{\,n-k}=2^{\,n/2+1}$,
\begin{equation}\label{eq:even-bound}
1-\rho(n)\ \le\ 2\sum_{g\in\RM(k,n)\setminus\{0\}}2^{-\mathrm{wt}(g)},
\end{equation}
and the contribution of the minimum-weight codewords,
\[
A_{w_0}\,2^{-w_0},\qquad
A_{w_0}=2^{\,k}\prod_{i=0}^{n-k-1}\frac{2^{\,n-i}-1}{2^{\,n-k-i}-1},
\]
satisfies $\log_2\bigl(A_{w_0}2^{-w_0}\bigr)\le \tfrac{n^2}{2}-2^{\,n/2+1}
\to-\infty$. In particular $A_{w_0}2^{-w_0}\to0$ super-exponentially.
\end{proposition}

\begin{theorem}[Asymptotics]\label{thm:asy}
\begin{enumerate}
\item (Even $n$.) Unconditionally, for every even $n$,
      $1-\rho(n)\le 2\sum_{g\ne0}2^{-\mathrm{wt}(g)}$
      (a finite quantity, \eqref{eq:even-bound}), whose minimum--weight
      part $2A_{w_0}2^{-w_0}$ tends to $0$ super-exponentially
      (Proposition~\ref{prop:evenmin}). Moreover
      $\lim_{n\to\infty,\ n\,\mathrm{even}}\rho(n)=1$ holds
      \emph{conditionally} on the classical Reed--Muller
      weight--enumerator input \eqref{eq:cited} below.
\item (Odd $n$.) For all odd $n\ge3$, $\;0<\rho(n)<1$ (both inequalities
      proved unconditionally; see Proposition~\ref{prop:oddbounds}). The
      value of $\lim_{n\to\infty,\ n\,\mathrm{odd}}\rho(n)$ is left
      \emph{open}; see Remark~\ref{rem:oddlimit} for the heuristic value
      $\prod_{i\ge1}(1-2^{-i})$ and the precise obstruction to proving it.
\end{enumerate}
Thus $\rho(n)$ does not converge: it splits by parity, with the even
subsequence tending to $1$ (conditionally on \eqref{eq:cited}) while the
odd subsequence stays strictly inside $(0,1)$.
\end{theorem}

\section{The reduction (Theorem~\ref{thm:reduction})}

\begin{proof}
A nonzero $g$ with $\deg g\le k$ is $g=\sum_{\mathrm{wt}(u)\le k}c_ux^u$
with coefficient vector $c=(c_u)\ne0$ indexed by the columns of $E$; its
value vector $(g(x))_x$ equals $E\,c$. Hence $g$ vanishes on $S$ iff
$E_S\,c=0$, and a nonzero such $c$ exists iff the column rank of $E_S$ is
$<M$. Now $g$ annihilates $f$ iff $g$ vanishes on $\supp(f)$, and $g$
annihilates $f\oplus1$ iff $g$ vanishes on $\F_2^n\setminus\supp(f)$. Thus
a nonzero annihilator of degree $\le k=m-1$ of $f$ or of $f\oplus1$ exists
iff $E_{\supp(f)}$ or $E_{\F_2^n\setminus\supp(f)}$ is column-rank
deficient. Since $\AI(f)\le m$ always, $\AI(f)=m$ iff no such annihilator
exists, i.e. iff both matrices have full column rank $M$.
\end{proof}

\begin{remark}\label{rem:fullrank}
The full matrix $E$ has column rank exactly $M$: its columns are the
evaluation vectors of the distinct monomials $x^u$, $\mathrm{wt}(u)\le k$,
and these are linearly independent over $\F_2$ because the algebraic normal
form of a Boolean function is unique (a nontrivial $\F_2$-combination of
monomials is a nonzero polynomial, hence a nonzero function). We use this
in Proposition~\ref{prop:oddbounds}.
\end{remark}

\section{The parity dichotomy (Theorem~\ref{thm:dichotomy})}

\begin{proof}
By $\binom nj=\binom n{n-j}$ and $\sum_{j=0}^n\binom nj=N$:

If $n$ is even, $k=n/2-1$ and the central term $\binom n{n/2}$ is left out
symmetrically, so $N=2M+\binom n{n/2}$, i.e.
$M=\tfrac12\bigl(N-\binom n{n/2}\bigr)=N/2-\tfrac12\binom n{n/2}<N/2$.

If $n$ is odd, $k=(n-1)/2$ and there is no central term, so the two halves
are equal: $N=2M$, i.e. $M=N/2$.
\end{proof}

\section{Odd $n$: complement symmetry (Theorem~\ref{thm:odd-sym})}
\label{sec:odd-sym}

Let $n$ be odd, $k=(n-1)/2$. View $\F_2^{N}$ as the space of functions
$\F_2^n\to\F_2$ with inner product
$\langle a,b\rangle=\sum_{x}a(x)b(x)$, which is nondegenerate, so
$\dim X+\dim X^{\perp}=N$ and $(X+Y)^{\perp}=X^{\perp}\cap Y^{\perp}$ for
all subspaces $X,Y$. Put
\[
V:=\RM(k,n)=\{g:\deg g\le k\},\qquad \dim V=M=N/2
\]
(by Theorem~\ref{thm:dichotomy}(2)), and for $T\subseteq\F_2^n$ let
$W_T:=\{a:\supp(a)\subseteq T\}$, $\dim W_T=|T|$.

\begin{lemma}\label{lem:selfdual}
$V^{\perp}=V$ (self-duality).
\end{lemma}
\begin{proof}
Reed--Muller duality: $\RM(r,n)^{\perp}=\RM(n-r-1,n)$. With $r=k=(n-1)/2$,
$n-r-1=(n-1)/2=k$, hence $V^{\perp}=\RM(k,n)=V$.
\end{proof}

\begin{lemma}\label{lem:WTdual}
$W_T^{\perp}=W_{T^c}$ for every $T$.
\end{lemma}
\begin{proof}
$a\perp W_T$ iff $\langle a,e_{x_0}\rangle=a(x_0)=0$ for every point
$x_0\in T$ (testing against indicators $e_{x_0}$, which span $W_T$), i.e.
iff $\supp(a)\subseteq T^c$.
\end{proof}

\begin{proof}[Proof of Theorem~\ref{thm:odd-sym}]
By the proof of Theorem~\ref{thm:reduction}, the right kernel of $E_S$ is
isomorphic to $\{g\in V:g|_S=0\}=V\cap W_{S^c}$. As $E_S$ is square of size
$M=N/2$,
\begin{equation}\label{eq:ker}
\rank E_S=M\iff V\cap W_{S^c}=\{0\}.
\end{equation}
We prove the stronger
$\dim(V\cap W_S)=\dim(V\cap W_{S^c})$. Using $\dim V=N/2$ and
$\dim W_S=N/2$,
\[
\dim(V\cap W_S)=\dim V+\dim W_S-\dim(V+W_S)=N-\dim(V+W_S).
\]
By nondegeneracy and Lemmas~\ref{lem:selfdual}--\ref{lem:WTdual},
\[
\dim(V+W_S)=N-\dim\bigl((V+W_S)^{\perp}\bigr)
=N-\dim\bigl(V^{\perp}\cap W_S^{\perp}\bigr)
=N-\dim\bigl(V\cap W_{S^c}\bigr).
\]
Substituting, $\dim(V\cap W_S)=\dim(V\cap W_{S^c})$. Hence one is $\{0\}$
iff the other is, and by \eqref{eq:ker} applied to $S$ and $S^c$,
$\rank E_S=M\iff\rank E_{S^c}=M$. By Theorem~\ref{thm:reduction}, for odd
$n$ membership in $B_n^*$ requires both $E_{\supp(f)}$ and
$E_{\supp(f)^c}$ invertible, conditions now seen to be equivalent; this
gives \eqref{eq:odd-rho}.
\end{proof}

\section{Even $n$: optimality is generic}\label{sec:even}

Let $n$ be even, $k=n/2-1$, $V=\RM(k,n)$, $w_0=2^{\,n-k}=2^{\,n/2+1}$.

\begin{lemma}\label{lem:ratio}
For $0\le w\le N/2$,
$\dfrac{\binom{N-w}{N/2}}{\binom{N}{N/2}}
=\prod_{i=0}^{w-1}\dfrac{N/2-i}{N-i}\le 2^{-w}$.
\end{lemma}
\begin{proof}
Direct cancellation gives the product; each factor
$\frac{N/2-i}{N-i}\le\frac12$ since $N/2-i\le\frac12(N-i)\iff i\ge0$.
\end{proof}

\begin{proof}[Proof of \eqref{eq:even-bound}]
Pick $f\in B_n$ uniformly, so $S=\supp(f)$ is a uniform $(N/2)$-subset.
By Theorem~\ref{thm:reduction} and the union bound (both $S,S^c$ uniform),
\[
1-\rho(n)\le 2\Pr[\rank E_S<M]
=2\Pr\bigl[\exists\,g\in V\setminus\{0\}:S\subseteq Z(g)\bigr]
\le 2\!\!\sum_{g\in V\setminus\{0\}}\!\!\Pr[S\subseteq Z(g)].
\]
For $g$ with $\mathrm{wt}(g)=w$, $|Z(g)|=N-w$ and
$\Pr[S\subseteq Z(g)]=\binom{N-w}{N/2}/\binom{N}{N/2}\le2^{-w}$ by
Lemma~\ref{lem:ratio}. This is \eqref{eq:even-bound}.
\end{proof}

\begin{lemma}[Minimum weight of $\RM(k,n)$]\label{lem:minw}
Every nonzero $g\in\RM(k,n)$ has $\mathrm{wt}(g)\ge w_0=2^{\,n-k}$
(minimum distance of $\RM(k,n)$), and the number of weight-$w_0$ codewords
is
\[
A_{w_0}=2^{\,k}\prod_{i=0}^{n-k-1}\frac{2^{\,n-i}-1}{2^{\,n-k-i}-1},
\qquad \log_2 A_{w_0}\ \le\ k+(k+1)(n-k)\ \le\ \tfrac{n^2}{2}.
\]
\end{lemma}
\begin{proof}
The minimum distance of $\RM(r,n)$ is $2^{\,n-r}$ and the number of
minimum-weight codewords is
$2^{\,r}\prod_{i=0}^{n-r-1}\frac{2^{\,n-i}-1}{2^{\,n-r-i}-1}$; take $r=k$.
We bound each of the $n-k$ factors. For integers $0\le i\le n-k-1$ we have
$n-k-i\ge1$, hence $2^{\,n-k-i}-1\ge 2^{\,n-k-i-1}$, and therefore
\[
\frac{2^{\,n-i}-1}{2^{\,n-k-i}-1}
<\frac{2^{\,n-i}}{2^{\,n-k-i-1}}=2^{\,k+1}.
\]
Consequently $A_{w_0}<2^{\,k}\cdot(2^{\,k+1})^{\,n-k}
=2^{\,k+(k+1)(n-k)}$, i.e. $\log_2 A_{w_0}\le k+(k+1)(n-k)$. With
$k=n/2-1$ this exponent equals $\tfrac{n^2}{4}+n-1$, and
$\tfrac{n^2}{2}-\bigl(\tfrac{n^2}{4}+n-1\bigr)=\bigl(\tfrac n2-1\bigr)^2\ge0$,
so $\log_2 A_{w_0}\le\tfrac{n^2}{2}$.
\end{proof}

\begin{proof}[Proof of Proposition~\ref{prop:evenmin}]
Inequality \eqref{eq:even-bound} is proved above. By Lemma~\ref{lem:minw}
the smallest weight is $w_0$, so the minimum-weight term in
\eqref{eq:even-bound} is $A_{w_0}2^{-w_0}$, and
$\log_2(A_{w_0}2^{-w_0})\le\frac{n^2}{2}-w_0=\frac{n^2}{2}-2^{\,n/2+1}\to
-\infty$.
\end{proof}

\paragraph{What is unconditional, and the corrected dominant--term values.}
For every even $n$, \eqref{eq:even-bound} is a rigorous, finite upper bound
for $1-\rho(n)$, and its minimum--weight part is exactly
\[
2A_{w_0}2^{-w_0}\in\{\,0.234\ (n=4),\ 0.0795\ (n=6),\
3.62\cdot10^{-4}\ (n=8),\ 9.32\cdot10^{-11}\ (n=10),\ \dots\,\},
\]
computed from the closed form for $A_{w_0}$; these numbers tend to $0$
super-exponentially by Proposition~\ref{prop:evenmin}. We stress that
$2A_{w_0}2^{-w_0}$ is \emph{not} itself an upper bound for $1-\rho(n)$:
the valid upper bound is the full sum in \eqref{eq:even-bound}, which also
contains positive contributions from codewords of weight $>w_0$. For
small $n$ those higher--weight contributions are comparable to the
dominant term: e.g.\ for $n=6$ the full sum in \eqref{eq:even-bound}
equals $0.1217$, of which $0.0795$ comes from the minimum--weight
codewords and $0.0422$ from the rest (these values were verified by direct
enumeration of $\RM(n/2-1,n)$).

\paragraph{Completion of Theorem~\ref{thm:asy}(1): the precise obstruction.}
By \eqref{eq:even-bound},
\begin{equation}\label{eq:tailsplit}
1-\rho(n)\ \le\ 2\sum_{w\ge w_0}A_w\,2^{-w}
\ =\ 2A_{w_0}2^{-w_0}\ +\ 2\sum_{w>w_0}A_w\,2^{-w},
\end{equation}
where $A_w$ is the number of weight-$w$ codewords of $\RM(k,n)$. The
minimum--weight term tends to $0$ super-exponentially
(Proposition~\ref{prop:evenmin}). We record an exact identity and then
isolate honestly the single analytic input needed for the tail.

\emph{Generating identity.} Let $R\subseteq\F_2^n$ be a random set in which
each point is included independently with probability $\tfrac12$. For a
fixed $g$ of weight $w$, $\Pr[g|_R=0]=2^{-w}$, so by linearity of
expectation
\begin{equation}\label{eq:genid}
\sum_{g\in\RM(k,n)}2^{-\mathrm{wt}(g)}
=\mathbb E_R\bigl[\#\{g\in\RM(k,n):g|_R=0\}\bigr]
=\mathbb E_R\bigl[2^{\,M-\rank E_R}\bigr].
\end{equation}
Thus the union sum in \eqref{eq:even-bound} equals
$\mathbb E_R\bigl[2^{M-\rank E_R}\bigr]-1$. (Identity \eqref{eq:genid}
holds for every $n$ and was verified numerically; it shows the union sum
is finite and $\ge0$.)

\emph{Why the tail is genuinely hard.} One might hope to bound
$\sum_{w>w_0}A_w2^{-w}$ crudely by splitting at $2w_0$ and using the
trivial estimate $A_w\le|V|=2^{M}$ for $w\ge 2w_0$. This \emph{fails}: for
even $n$ one has $M=N/2-\tfrac12\binom n{n/2}$ while $2w_0=2^{\,n/2+2}$,
and a direct check shows $M-2w_0\to+\infty$ (e.g.\ $M-2w_0=29$ at $n=8$
and $\approx1.1\cdot10^{8}$ at $n=28$), so the crude tail bound
$2^{M}\cdot2^{-2w_0}$ \emph{diverges}. Controlling the tail therefore
requires the genuinely sharp Reed--Muller weight--distribution bounds, not
a trivial counting estimate.

\emph{Classical input.} Consequently the completion of part (1) rests on
the following statement about the weight distribution of the low--order
Reed--Muller code $V=\RM(n/2-1,n)$:
\begin{equation}\label{eq:cited}
\sum_{g\in V\setminus\{0\}}2^{-\mathrm{wt}(g)}\ \longrightarrow\ 0
\qquad(n\to\infty,\ n\ \text{even}).
\end{equation}
This is a consequence of the now--standard sharp bounds on the weight
distribution of Reed--Muller codes (Kasami--Tokura for weights in
$[w_0,2w_0)$, and the modern bounds of Sberlo--Shpilka and of Samorodnitsky
for the intermediate range, combined with the exponential factor $2^{-w}$
that controls the bulk near $w=N/2$). We do \emph{not} reprove these
Reed--Muller weight bounds here; \eqref{eq:cited} is the single external
ingredient of part (1), and we make no claim to have established it
elementarily. We also emphasise the correction of an earlier, false
formulation: it is \emph{not} true that
$\sum_{w>w_0}A_w2^{-w}=o\bigl(A_{w_0}2^{-w_0}\bigr)$ for small $n$ (for
$n=6$ the tail $0.0422$ is of the same order as the dominant $0.0795$);
the correct sufficient statement is the convergence \eqref{eq:cited} of the
\emph{whole} union sum.

Granting \eqref{eq:cited}, \eqref{eq:even-bound} gives
$1-\rho(n)\le 2\sum_{g\ne0}2^{-\mathrm{wt}(g)}\to0$, hence $\rho(n)\to1$
for even $n$.

\emph{Unconditional numerical check of \eqref{eq:even-bound}.} Direct
computation of the weight enumerator of $\RM(n/2-1,n)$ gives
\[
1-\rho(4)\le 0.234,\quad 1-\rho(6)\le 0.122,\quad 1-\rho(8)\le 2.0\cdot10^{-3},
\]
(the last computed via \eqref{eq:genid}), and direct Monte-Carlo over
balanced $f$ gives $1-\rho(8)\approx 3\cdot10^{-4}$, comfortably below the
bound.

\section{Odd $n$: bounds and the open limit}\label{sec:odd-asy}

\begin{proposition}\label{prop:oddbounds}
For every odd $n\ge3$, $\;0<\rho(n)<1$.
\end{proposition}
\begin{proof}
\emph{$\rho(n)>0$ (self-contained).} For odd $n$, $E$ is an $N\times M$
matrix with $M=N/2$ and column rank $M$ (Remark~\ref{rem:fullrank}).
Hence $E$ has $M$ linearly independent rows; let $S\subseteq\F_2^n$ be
their index set, so $|S|=M=N/2$ and the square submatrix $E_S$ is
invertible. The balanced function $f=\mathbf 1_S$ then has
$\rank E_S=M$, and by Theorem~\ref{thm:odd-sym} also $\rank E_{S^c}=M$;
by Theorem~\ref{thm:reduction}, $f\in B_n^*$. Thus $B_n^*\ne\varnothing$
and $\rho(n)>0$. (This reproves, with an elementary argument, the
classical fact that balanced functions of optimal algebraic immunity
exist for every $n$.)

\emph{$\rho(n)<1$.} Let $k=(n-1)/2$ and let $H\subseteq\F_2^n$ be an
affine subspace (flat) of codimension $k$, so $|H|=2^{\,n-k}=2^{(n+1)/2}=w_0$.
The indicator $\mathbf 1_H$ is a Boolean function of degree exactly $k$
(the indicator of a codimension-$k$ flat is a product of $k$ affine forms,
hence has degree $k$); thus $\mathbf 1_H\in V\setminus\{0\}$ is a
nonzero codeword of $\RM(k,n)$. Since $w_0=2^{(n+1)/2}\le 2^{\,n-1}=N/2$
for $n\ge3$, the complement $\F_2^n\setminus H$ has $N-w_0\ge N/2$ points,
so there exist $(N/2)$-subsets $S\subseteq\F_2^n\setminus H$. For any such
$S$, $\mathbf 1_H$ vanishes on $S$, so $E_S$ is singular and the balanced
$f=\mathbf 1_S\notin B_n^*$. Quantitatively, for a fixed flat $H$,
\[
\Pr[S\subseteq\F_2^n\setminus H]
=\frac{\binom{N-w_0}{N/2}}{\binom{N}{N/2}}
=\prod_{i=0}^{w_0-1}\frac{N/2-i}{N-i}\ >\ 0,
\]
so $1-\rho(n)>0$, i.e. $\rho(n)<1$, for every odd $n\ge3$.
\end{proof}

\begin{remark}[On a uniform lower bound for $1-\rho(n)$]\label{rem:notuniform}
We do \emph{not} claim that $\rho(n)$ is bounded away from $1$ as
$n\to\infty$, and we caution that the elementary one--codeword estimate
above does \emph{not} prove it: the bound
$\prod_{i=0}^{w_0-1}\frac{N/2-i}{N-i}$ is at most $2^{-w_0}\to0$
super-exponentially. A second--moment attempt over the family of
minimum--weight codewords also fails to give a uniform bound, since the
expected number of minimum--weight flats vanishing on a random $S$,
namely $A_{w_0}\binom{N-w_0}{N/2}/\binom{N}{N/2}$, already tends to $0$
(e.g.\ it is $\approx4.4\cdot10^{-3}$ at $n=9$). Establishing a uniform
positive lower bound on $1-\rho(n)$ for odd $n$ is therefore left open; it
would follow from (and is essentially equivalent to) the odd limit being
$<1$, discussed next.
\end{remark}

\begin{remark}[The heuristic odd limit and the obstruction]\label{rem:oddlimit}
By Theorem~\ref{thm:odd-sym}, $\rho(n)$ is the probability that the square
$\tfrac N2\times\tfrac N2$ matrix $E_S$ over $\F_2$ is invertible, where
its $N/2$ rows are the degree-$\le k$ monomial evaluations at a uniformly
random $\tfrac N2$-subset of points. If those rows were independent uniform
vectors of $\F_2^{N/2}$, the invertibility probability would be exactly
$\prod_{i=1}^{N/2}(1-2^{-i})\to\prod_{i\ge1}(1-2^{-i})=0.288788\ldots$ as
$n\to\infty$. The rows are \emph{not} independent uniform (they are
Reed--Muller evaluation vectors, subject to the linear relations of the
self-dual code $V$), so this is only a heuristic.

The available numerical evidence is \emph{inconclusive}. Exact enumeration
gives $\rho(3)=0.800$. Monte-Carlo estimates give
$\rho(5)\approx0.32$, $\rho(7)\approx0.14$, and $\rho(9)\approx0.28$--$0.36$
(with substantial sampling variance at $n=9$). This sequence is strongly
non-monotone and the dip $\rho(7)\approx0.14$ lies well below
$0.2888$, so the data neither confirm nor refute the heuristic limit. We
therefore record
\[
\lim_{\substack{n\to\infty\\ n\ \mathrm{odd}}}\rho(n)
\stackrel{?}{=}\prod_{i\ge1}\bigl(1-2^{-i}\bigr)=0.288788\ldots
\]
only as a conjecture. Proving (or disproving) it is the genuinely hard
analytic core: it asks whether the algebraic dependence among
Reed--Muller evaluation rows leaves the limiting singularity probability of
$E_S$ unchanged from that of a Haar-random $\F_2$-matrix. No standard
theorem delivers this, and we leave it open.
\end{remark}

\section{Conclusion}

The exact characterisation $B_n^*=\{f\in B_n:\rank E_{\supp(f)}=\rank
E_{\supp(f)^c}=M\}$ (Theorem~\ref{thm:reduction}) together with the parity
dichotomy (Theorem~\ref{thm:dichotomy}) splits the problem completely by
the parity of $n$:
\begin{itemize}
\item \emph{Even $n$}: $E_S$ is tall ($M<N/2$), full column rank is
generic, and the dominant deviation $2A_{w_0}2^{-w_0}\to0$ is computed
unconditionally in Proposition~\ref{prop:evenmin}. The limit
$\rho(n)\to1$ follows from the classical Reed--Muller weight distribution
\eqref{eq:cited}; we showed that no elementary counting estimate suffices
for the tail, so this one input is essential.
\item \emph{Odd $n$}: $E_S$ is square, the complement symmetry
(Theorem~\ref{thm:odd-sym}) makes $\rho(n)$ a single
matrix-invertibility probability with $0<\rho(n)<1$ proved
unconditionally (Proposition~\ref{prop:oddbounds}). The value of the
odd limit is open; the heuristic $\prod_{i\ge1}(1-2^{-i})=0.288788\ldots$
is supported only by inconclusive, non-monotone numerics
(Remark~\ref{rem:oddlimit}).
\end{itemize}
Hence $\rho(n)$ does not converge; it splits by parity, with the even
subsequence tending to $1$ (conditionally on \eqref{eq:cited}) and the odd
subsequence remaining strictly inside $(0,1)$ with an unresolved limit.
\end{document}
